%%%%%%%%%%%%%%%%%%%%%%%%%%%%%%%%%%%%%%%%%
% Template de Relatório
% Version 0.8 (12/04/19)
%
% Authors:
% Franklin Coêlho
% Eudisley Anjos
% 
% License:
% CC BY-NC-SA 3.0 (http://creativecommons.org/licenses/by-nc-sa/3.0/)
%%%%%%%%%%%%%%%%%%%%%%%%%%%%%%%%%%%%%%%%%

\documentclass[10pt, a4paper, twocolumn]{article} 

\input{config.tex}



%%%%%%%%%%%%%%%%%%%%% HEADER %%%%%%%%%%%%%%%%%%%%%


\title{Título do Projeto}

\author{
    \centering\authorstyle Aluno 01\textsuperscript{1}, Aluno 02\textsuperscript{1}, Aluno 03\textsuperscript{1},..., Aluno n\textsuperscript{1}\\
    \textsuperscript{1}\institution{Bacharelando em Engenharia de Computação, Centro de Informática -- Universidade Federal da Paraíba}\\
    \email{aluno01@eng.ci.ufpb.br, aluno02@eng.ci.ufpb.br, alunon@ci.ufpb.br}
}

\date{}

%%%%%%%%%%%%%%%%%%%%%%%%%%%%%%%%%%%%%%%%%%%%%%%%%%%%%%

\begin{document}

\maketitle        

\thispagestyle{firstpage}

\lettrineabstract{Aqui começa o resumo. O Lorem Ipsum é um texto modelo da indústria tipográfica e de impressão. O Lorem Ipsum tem vindo a ser o texto padrão usado por estas indústrias desde o ano de 1500, quando uma misturou os caracteres de um texto para criar um espécime de livro. Este texto não só sobreviveu 5 séculos, mas também o salto para a tipografia electrónica, mantendo-se essencialmente inalterada. Foi popularizada nos anos 60 com a disponibilização das folhas de Letraset, que continham passagens com Lorem Ipsum, e mais recentemente com os programas de publicação como o Aldus PageMaker que incluem versões do Lorem Ipsum.} \\ \\ %Pula duas linhas
\textbf{Palavras-chave:} Palavra 1, Palavra 2, ... .

\section*{Lista de Siglas}
%Você pode ajustar a distância a partir do parâmetro "\hspace{}"
%Tente chegar em uma distância alinhada com as demais siglas
\textbf{UFPB} \hspace{1.5cm} Universidade Federal da Paraíba \\
\textbf{CI} \hspace{1.95cm} Centro de Informática \\
\textbf{VHDL} \hspace{1.5cm} Hardware Description Language \\
\textbf{FPGA} \hspace{1.5cm} Field Programmable Gate Array

\section*{INTRODUÇÃO}

Descreva os objetivos principais e secundários do seu projeto, contendo o que você espera obter. \par %Faz um parágrafo

O Lorem Ipsum é um texto modelo da indústria tipográfica e de impressão. O Lorem Ipsum tem vindo a ser o texto padrão usado por estas indústrias desde o ano de 1500, quando uma misturou os caracteres de um texto para criar um espécime de livro. Este texto não só sobreviveu 5 séculos, mas também o salto para a tipografia electrónica, mantendo-se essencialmente inalterada. Foi popularizada nos anos 60 com a disponibilização das folhas de Letraset, que continham passagens com Lorem Ipsum, e mais recentemente com os programas de publicação como o Aldus PageMaker que incluem versões do Lorem Ipsum.

%Aqui você pode inserir itens
\begin{itemize}
\item Item 1; 
\item Item 2; 
\item Item 3; 
\item ...
\end{itemize}

\section*{METODOLOGIA} 

Texto sobre a metodologia. \par

O Lorem Ipsum é um texto modelo da indústria tipográfica e de impressão. O Lorem Ipsum tem vindo a ser o texto padrão usado por estas indústrias desde o ano de 1500, quando uma misturou os caracteres de um texto para criar um espécime de livro. Este texto não só sobreviveu 5 séculos, mas também o salto para a tipografia eletrônica, mantendo-se essencialmente inalterada. Foi popularizada nos anos 60 com a disponibilização das folhas de Letraset, que continham passagens com Lorem Ipsum, e mais recentemente com os programas de publicação como o Aldus PageMaker que incluem versões do Lorem Ipsum.

\section*{DESCRIÇÃO DO PROJETO}

Descreva o funcionamento do projeto. \par

O Lorem Ipsum é um texto modelo da indústria tipográfica e de impressão. O Lorem Ipsum tem vindo a ser o texto padrão usado por estas indústrias desde o ano de 1500, quando uma misturou os caracteres de um texto para criar um espécime de livro. Este texto não só sobreviveu 5 séculos, mas também o salto para a tipografia eletrônica, mantendo-se essencialmente inalterada. Foi popularizada nos anos 60 com a disponibilização das folhas de Letraset, que continham passagens com Lorem Ipsum, e mais recentemente com os programas de publicação como o Aldus PageMaker que incluem versões do Lorem Ipsum.

%Inserção de imagem
\begin{figure}[H]
    \centering
    %"width" serve para ajustar a dimensão da imagem
    \includegraphics[width=0.4\linewidth]{pictures/logo_ci.png}
    \caption{Aqui você pode inserir a legenda da imagem.}
    \label{fig:f01}
\end{figure}

A Figura \ref{fig:f01} mostra a logomarca do CI. Você pode referenciar figuras, tabelas e as próprias referências ao final da proposta.
Para fazer uma referência, basta seguir assim: Segundo \citep{ackerman2007business}, podemos prosseguir com a proposta. Perceba que a referência automaticamente é gerada ao final do documento. Legal, não é?


%Aqui você pode inserir itens
\begin{itemize}
\item Item 1; 
\item Item 2; 
\item Item 3; 
\item ...
\end{itemize}


\section*{EXECUÇÃO DO PROJETO, TESTES E RESULTADOS}

Texto sobre a execução e testes do projeto. \par

O Lorem Ipsum é um texto modelo da indústria tipográfica e de impressão. O Lorem Ipsum tem vindo a ser o texto padrão usado por estas indústrias desde o ano de 1500, quando uma misturou os caracteres de um texto para criar um espécime de livro. Este texto não só sobreviveu 5 séculos, mas também o salto para a tipografia eletrônica, mantendo-se essencialmente inalterada. Foi popularizada nos anos 60 com a disponibilização das folhas de Letraset, que continham passagens com Lorem Ipsum, e mais recentemente com os programas de publicação como o Aldus PageMaker que incluem versões do Lorem Ipsum.

\begin{table}[h!]
\centering
\scriptsize
\begin{tabular}{|c|c|c|c|c|}
	\hline
    Atividade & \multicolumn{4}{c|}{Mês/Ano} \\
    \hline
    & dezembro/2018 & janeiro/2019 & fevereiro/2019 & março/2019  \\
    \hline
	1 & X & X & &   \\
    \hline
    2 & X & X & &   \\
    \hline
    3 &  &  & X &   \\
    \hline
    4 &  &  & & X   \\
    \hline
    5 &  &  & & X  \\
    \hline
\end{tabular}
\caption{Aqui temos a legenda da tabela.}
\label{tb-cronograma}
\end{table}

\textbf{Obs.: Há detalhes do roteiro de escrita no arquivo que fala dos projetos. Lembre-se de apresentar o código final do projeto.}

A seguir um exemplo de como inserir código no relatório: \\

% Definindo novas cores
\definecolor{verde}{rgb}{0,0.5,0}
% Configurando layout para mostrar codigos em C
\lstset{
  language=C,
  basicstyle=\ttfamily\small, 
  keywordstyle=\color{blue}, 
  stringstyle=\color{verde}, 
  commentstyle=\color{red}, 
  extendedchars=true, 
  showspaces=false, 
  showstringspaces=false, 
  numbers=left,
  numberstyle=\tiny,
  breaklines=true, 
  backgroundcolor=\color{gray!10},
  breakautoindent=true, 
  captionpos=b,
  xleftmargin=0pt,
}
\pagestyle{empty}
%Escrita do código
\begin{lstlisting}
#include <stdio.h>

int main(){
 
  int i;
  float dedicacaoCL1, media = 0;
  
  printf("Digite a sua dedicacao na disciplina Circuitos Logicos I: ");
  scanf("%f", &dedicacaoCL1);
  
  for (i = 0; i < dedicacaoCL1; i++){
     media++;
  }
   
   printf("Sua media em Circuitos Logicos I foi %.2f\n", media);
   /*Quanto maior a dedicacao, maior o resultado*/
   
   return 0;
}
\end{lstlisting}

\section*{CONCLUSÕES}

Texto sobre a conclusão. \par

O Lorem Ipsum é um texto modelo da indústria tipográfica e de impressão. O Lorem Ipsum tem vindo a ser o texto padrão usado por estas indústrias desde o ano de 1500, quando uma misturou os caracteres de um texto para criar um espécime de livro. Este texto não só sobreviveu 5 séculos, mas também o salto para a tipografia eletrônica, mantendo-se essencialmente inalterada. Foi popularizada nos anos 60 com a disponibilização das folhas de Letraset, que continham passagens com Lorem Ipsum, e mais recentemente com os programas de publicação como o Aldus PageMaker que incluem versões do Lorem Ipsum.

\textbf{Observação:} Lembrem de enviar junto ao relatório os arquivos fontes necessários para desenvolvimento do projeto. Estes podem incluir, arquivos do Quartus II, simulações, etc. 

%Você pode alterar as referências ao lado, na aba "refs.bib"
\printbibliography[title={REFERÊNCIAS}]

\end{document}